\documentclass[11pt,a4paper]{article}
\usepackage{amsmath}
\usepackage{amssymb}
\usepackage{amsthm}
\usepackage{latexsym}
\usepackage{pstricks}
\usepackage{amsbsy}
\usepackage{amscd}
\usepackage{mathrsfs}
\usepackage{tipa}
\usepackage{txfonts}
\usepackage{amsfonts}
\usepackage{graphicx}
\usepackage{listings}
\usepackage{amsmath}
\usepackage{amssymb}
\usepackage[all]{xy}
\usepackage{enumitem}
\setenumerate[1]{itemsep=1pt,partopsep=0pt,parsep=\parskip,topsep=6pt}
\setitemize[1]{itemsep=1pt,partopsep=0pt,parsep=\parskip,topsep=6pt}
\setdescription{itemsep=1pt,partopsep=0pt,parsep=\parskip,topsep=6pt}
\input xypic
\addtolength{\textheight}{1cm}
\addtolength{\textwidth}{3.8cm}
\addtolength{\oddsidemargin}{-1,8cm}
\addtolength{\evensidemargin}{-1,8cm}
\newtheorem{thm}{Theorem}[section]
\newtheorem{cor}[thm]{Corollary}
\newtheorem{lem}[thm]{Lemma}
\newtheorem{cond}{Condition}
\newtheorem{que}[thm]{Question}
\newtheorem{exm}[thm]{Example}
\newtheorem{prop}[thm]{Proposition}
\newtheorem{defn}[thm]{Definition}
\newtheorem{rem}[thm]{Remark}
\newtheorem{nota}[thm]{Notation}


\parindent15pt

\begin{document}


\title{\Large \bf  On  a  proof of  the Gorenstein  Symmetry Conjecture
\footnotetext{
*Corresponding author.  E-mail address: sdwg001@163.com.\\
This research was supported by the Natural Science Foundation of Hunan Province (No. 2023JJ30561).
 }}



 \vskip 0.5cm
\author{{\normalsize  Weiqing Li }\\
{\small College of Mathematics and Statistics, Jishou University, Jishou 416000,  P. R. China}\date{}}
 \maketitle
 \vskip-1cm


\begin{abstract}  We give  a homological proof of  the famous Gorenstein Symmetry Conjecture, which asserts that   one-sided finiteness of the self-injective
dimension of an Artin algebra is sufficient for Gorensteinness.
\end{abstract}







\def\s{\stackrel}
\def\Longrightarrow{{\longrightarrow}}
\def\A{\mathcal{A}}
\def\B{\mathcal{B}}
\def\C{\mathcal{C}}
\def\D{\mathcal{D}}
\def\T{\mathcal{T}}
\def\R{\mathcal{R}}
\def\P{\mathcal{P}}
\def\S{\mathcal{S}}
\def\H{\mathcal{H}}
\def\U{\mathscr{U}}
\def\V{\mathscr{V}}
\def\M{\mathscr{M}}
\def\N{\mathcal{N}}
\def\W{\mathscr{W}}
\def\X{\mathscr{X}}
\def\Y{\mathscr{Y}}
\def\Z{\mathcal {Z}}
\def\I{\mathcal {I}}
\def\add{\mbox{add}}
\def\Aut{\mbox{Aut}}
\def\coker{\mbox{coker}}
\def\deg{\mbox{deg}}
\def\diag{\mbox{diag}}
\def\dim{\mbox{dim}}
\def\End{\mbox{End}}
\def\Ext{\mbox{Ext}}
\def\Hom{\mbox{Hom}}
\def\Gr{\mbox{Gr}}
\def\id{\mbox{id}}
\def\Im{\mbox{Im}}
\def\ind{\mbox{ind}}
\def\mod{\mbox{mod}}
\def\mul{\multiput}
\def\c{\circ}
\def \text{\mbox}

\hyphenation{ap-pro-xi-ma-tion}



\textbf{Key words:}    Gorenstein Symmetry Conjecture;   Artinian rings;  Cotorsion theory

\medskip

\textbf{2020 Mathematics Subject Classification:}  16G10; 18G25.



\section{ Introduction}



Throughout this paper, $\Lambda$ is an associative ring with identity and
all $\Lambda$-modules are unitary  right  $\Lambda$-modules. We view left $\Lambda$-modules as right  ones over
the opposite ring $\Lambda^{{\rm op}}$. We denote by (mod-$\Lambda$)  Mod-$\Lambda$ the class of all  (finitely generated) right    $\Lambda$-modules.
The symbols  id$_{\Lambda}$$(M)$ and pd$_{\Lambda}$$(M)$  stand  for the   injective and projective   dimension of
  a  right $\Lambda$-module $M$, respectively.

  The main purpose of this paper is to give a positive
answer to the following  long-standing conjecture posed by  Auslander and Reiten in \cite[p. 150]{ar1} (see also \cite[Conjecture (13)]{ars}).
\\
\\
{\bf Gorenstein Symmetry Conjecture (GSC) }\ \ {\it  Let   $\Lambda$  be    an Artin algebra   with {\rm id}$_{\Lambda}$$(\Lambda) < \infty$. Then {\rm id}$_{\Lambda^{{\rm op}}}$$(\Lambda^{{\rm op}}) < \infty$}.
\\
\par
For an Artin algebra $\Lambda$,   we learn from \cite[p. 121]{ar1}  and  \cite[Corollary 2]{hua} that {\rm id}$_{\Lambda}$$(\Lambda) \leq 1$ if and only if {\rm id}$_{\Lambda^{{\rm op}}}$$(\Lambda^{{\rm op}}) \leq 1$.
Later, Beligiannis \cite[Theorem 11.4]{bel}  proved that GSC holds for the class of  virtually Gorenstein Artin  algebras.

We show in Theorem \ref{thm4}   that GSC is true for any Artin  algebra.
\\
\par
We continue the introduction by recalling some basic definitions and notions
that we use in this paper. Given a class $ \mathscr{C}$ of    $\Lambda$-modules, we denote

  \begin{itemize}
 \item  $^\perp {\mathscr{C}}=\{X: {\rm Ext_{\Lambda}^{1}}(X, C)=0$  for all $C\in
{\mathscr {C}}\}$,
 \item  $ {\mathscr {C}}^\perp=\{X: {\rm Ext_{\Lambda}^{1}}(C, X)=0$
for all $C\in {\mathscr {C}}\}$,
\item  $^{\perp_{\infty}} {\mathscr{C}}=\{X: {\rm Ext}_{\Lambda}^{i}(X, C)=0$  for all $C\in
{\mathscr {C}} $  and $i \geq 1$\},
 \item  $ {\mathscr {C}}^{\perp_{\infty}}=\{X: {\rm Ext}_{\Lambda}^{i}(C, X)=0 $
for all $C\in {\mathscr {C}}$  and $i \geq 1$\},
\item${\bf Proj}(\Lambda)$
= the class of all projective right $\Lambda$-modules,
\item${\bf Inj}(\Lambda)$
= the class of all injective right $\Lambda$-modules.\end{itemize}


{\bf   (Pre)covers and (pre)envelopes.} \ \ Let $\mathscr{C}$    be a  class of modules in  Mod-$\Lambda$.
A  morphism $\phi: C\rightarrow M$ of Mod-$\Lambda$ with $C\in \mathscr{C}$ is called a  $\mathscr{C}$-$precover$ \cite{eno5} of $M$ if for any
morphism $f: C'\rightarrow M$ with $C'\in \mathscr{C}$,
there is a morphism $g : C'\rightarrow C$ such that $\phi g = f$. Moreover, if the only such $g$ are automorphisms of $C$ when
$C'= C$ and $f= \phi $, then the $\mathscr{C}$-precover $\phi$ is
called a $ \mathscr{C}$-$cover$. Dually, we have the
definitions of a  $\mathscr{C}$-preenvelope and a $\mathscr{C}$-envelope. We say that $\mathscr{C}$ is $(pre)covering$
($(pre)enveloping)$ in case every module in  Mod-$\Lambda$  has a $\mathscr{C}$-(pre)cover ($\mathscr{C}$-(pre)envelope).

A  module $M$ is said to have a  {\it special}  $\mathscr{C}$-$precover$, if there is an  exact sequence  $ 0
\rightarrow K\rightarrow C\rightarrow M\rightarrow 0$  with $C\in
\mathscr{C}$ and $K\in
{\mathscr {C}}^\perp$.  $M$ is said to have a  {\it special}  $\mathscr{C}$-$preenvelope$, if there is an  exact sequence  $ 0
\rightarrow M\rightarrow C\rightarrow L\rightarrow 0$  with $C\in
\mathscr{C}$ and $L\in
{^\perp {\mathscr{C}}}$.
\\
\par
{\bf   Cotorsion  theory.}\ \ A pair $(\mathscr{X},\mathscr{Y})$  of subcategories of  Mod-$\Lambda$   is said to be  a  {\it cotorsion  theory} \cite{eno5} if ${\mathscr{X}}^\perp = {\mathscr{Y}}
 $ and $^\perp{\mathscr {Y}}$ = ${\mathscr {X}}$. A cotorsion theory (${\mathscr{X},\ \mathscr{Y}}$)
 is called $complete$   if every object in  Mod-$\Lambda$   has a   special  ${\mathscr{X}}$-precover
and  a   special  $\mathscr
{Y}$-preenvelope.  A
cotorsion theory (${\mathscr{X},\ \mathscr{Y}}$)
 is said to be  $ hereditary$   if whenever
 $0\rightarrow Y\rightarrow Y'\rightarrow
Y''\rightarrow 0$ is exact with $Y$, $Y'\in {\mathscr{Y}}$
then $Y''\in {\mathscr{Y}}$.

\medskip

\section{ A characterization of    modules of  virtually finite  injective
  dimension}


In this section, we   give a characterization of  {\it modules of  virtually finite  injective
  dimension} \cite{bel} over  right Noetherian rings.

First, we recall the definitions of Gorenstein projective and   injective
   modules introduced by Enochs and Jenda in \cite{eno9}.


\begin{defn} \label{defn1} \rm{
 A  {\it  complete projective resolution} is an exact sequence of projective
 $\Lambda$-modules,  $$\cdots  \rightarrow P_{1} \rightarrow P_{0}\rightarrow
P_{-1}\rightarrow P_{-2}\rightarrow \cdots,  $$
 such that Hom$_{\Lambda}$($-$, $Q$) leaves the sequence exact whenever $Q$ is a projective   $\Lambda$-module; the module $M$ = im($P_{0}\rightarrow
P_{-1}$) is then  said to be {\it Gorenstein  projective}.

{\it Gorenstein  injective }  modules are defined dually.

 We use $\mathscr{GP}$ and $\mathscr{GI}$    to denote, respectively,  the classes of all Gorenstein  projective and  injective
   (right)  $\Lambda$-modules.  The modules in $^\perp\mathscr{GI}$ are called {\it modules of  virtually finite  injective
  dimension} \cite{bel}.

We need   the following    proposition  due to $\check{\rm S}$aroch and $\check{\rm S}$$\check{\rm t}$ov$\acute{\rm {\i}}$$\check{\rm c}$ek.}\end{defn}
\begin{prop}\label{prop1} {\rm (\cite[Theorem 5.6]{ss})}. For any ring $\Lambda$, $(^\perp\mathscr{GI}$, $\mathscr{GI})$ is a hereditary complete
  cotorsion theory.   \end{prop}

We recall some definitions from \cite{gbe} and \cite{eno5}.

 \begin{defn} \label{defn2} \rm{Let $\mu$ be an ordinal and $\mathscr{A}$ = ($A_{\alpha} $ $|$ $\alpha \leq \mu$) be a sequence of modules.
  If $A_{0} $ = 0,  $A_{\alpha} \subseteq  A_{\alpha + 1}$ for all $\alpha < \mu$ and   $A_{\alpha} $ = $\bigcup_{\beta < \alpha} A_{\beta} $ for all
 limit ordinals $\alpha \leq \mu$, then the sequence $\mathscr{A}$  is called a {\it continuous {\rm (}well ordered{\rm )} chain of modules.}

   Let $M$ be a module and $\mathscr{X}$ be a class of modules. $M$ is  $\mathscr{X}$-{\it filtered}, provided that there is an ordinal $\mu$
   and  a continuous chain of modules, ($M_{\alpha} $ $|$ $\alpha \leq \mu$), consisting of submodules of $M$ such that $M$ = $M_{\mu} $,
    and each of the modules $M_{\alpha + 1}/M_{\alpha}  $  ($\alpha < \mu$) is isomorphic to an element of $\mathscr{X}$. Then chain ($M_{\alpha} $ $|$ $\alpha \leq \mu$)
   is called an $\mathscr{X}$-{\it filtration} of $M$. If $\mu$ if finite, then  $M$ is said to be {\it finitely} $\mathscr{X}$-{\it filtered}.}\end{defn}

For convenience, we list the following classical result from \cite{gbe}.

\begin{lem}\label{lem2}  {\rm (\cite[Corollary 6.14]{gbe})}. Let   $\mathscr{X}$ be a set of modules containing $\Lambda$.
Then the class  $^\perp(\mathscr{X}^\perp)$ consists of all direct summands of $\mathscr{X}$-filtered modules.
\end{lem}


We get a useful  corollary.

\begin{cor}\label{cor1}   Let   $\mathscr{X}$ be a set of modules containing $\Lambda$.
Then  any   simple module in    $^\perp(\mathscr{X}^\perp)$ is  isomorphic to a  direct summand of some module in $\mathscr{X}$. \end{cor}

\begin{proof} Let  $S$  be a   simple module in    $^\perp(\mathscr{X}^\perp)$. Then $S$    is isomorphic to a  direct summand of $M$ which is  $\mathscr{X}$-{\it filtered} by
Lemma  \ref{lem2}. There is an ordinal $\mu$
   and  a continuous (well ordered) chain ($M_{\alpha} $ $|$ $\alpha \leq \mu$) of submodules  of $M$ such that $M$ = $M_{\mu} $,
    and each of the modules $M_{\alpha + 1}/M_{\alpha}  $  ($\alpha < \mu$) is isomorphic to an element of $\mathscr{X}$.
    We may assume that $S$  is a  submodule  of $M$. Then the subset
 \\
 \\
\centerline{   $\kappa$   = \{$\alpha \leq \mu$  \ $|$\  $S \subseteq M_{\alpha}$\}}
 \\
 \\
 is nonempty. Note that $\mu$ is well ordered. Hence  $\kappa$  has a least element (cf. \cite[Definition 1.1.4]{eno5}), say $\beta$.
 If  $\beta$ is a limit ordinal, then   $M_{\beta} = \bigcup_{\gamma < \beta} M_{\gamma} $. But $S   $ is simple. Thus $S \subseteq M_{\gamma}$ for some
 $\gamma < \beta$, in contradiction to minimality. Thus $\beta$ is not a limit ordinal. Since  $M_{0} = 0$, we have that $ \beta \neq 0$. Therefore $\beta$ is a successor ordinal,
 i.e., $\beta$ = $\lambda + 1$ for some $\lambda < \mu$. Then $S  $ is not contained in  $  M_{\lambda}$, and hence $S \cap M_{\lambda} = \{0\}$ since  $S  $ is simple.
Note that $S$    is   a  direct summand of $M$. So $S$    is also   a  direct summand of $ M_{\lambda + 1}$. It follows that   $S$    is isomorphic to
   a  direct summand of  $M_{\lambda + 1}/M_{\lambda}  $ which is isomorphic to an element of $\mathscr{X}$. This finishes the proof. \end{proof}

The following homological lemma is a generalization of \cite[Theorem 7.3.4]{eno5}.

\begin{lem}\label{lem1}  {\rm (\cite[Lemma 3.8]{mao}.)} Let $M$ and $N$ be   $\Lambda$-modules,    $n$  a positive integer and $M$
the union of a continuous chain of submodules $(M_{\alpha} $ $|$ $\alpha \leq \kappa)$. If ${\rm Ext}_{\Lambda}^{n}(M_{0}, N)=0$ and ${\rm Ext}_{\Lambda}^{n}(M_{\alpha + 1}/M_{\alpha}, N)=0$
whenever  $\alpha + 1 < \kappa$\, then ${\rm Ext}_{\Lambda}^{n}(M, N)=0$ \end{lem}

To simplify the proof  of the next theorem, we give the following notation.
\begin{nota}\label{nota1}  {\rm  \  Let  $\Lambda$ be a right  Noetherian ring. Then there is a
family ($E_{j})_{j\in J}$ of indecomposable injective right  $\Lambda$-modules such that every injective right  $\Lambda$-module is the direct sum
of copies of the various $E_{j}$ (see the proof of \cite[Theorem 5.4.1]{eno5}). For each $E_{j}$, take a  projective resolution $$ \cdots \stackrel {f_{j2}} \longrightarrow
{F_{j1}}\stackrel {f_{j1}} \longrightarrow {F_{j0}}\stackrel {f_{j0}}
\longrightarrow E_{j}  \longrightarrow  0  \eqno (\spadesuit_{j})$$
 of  $E_{j}$, and let
 $$ 0  \longrightarrow \Lambda_{\Lambda}\stackrel {h_{0}}
\longrightarrow
E_0\stackrel {h_{1}} \longrightarrow E_1\stackrel {h_{2}}
\longrightarrow \cdots \eqno (\clubsuit)$$
  be a minimal  injective resolution of   $\Lambda_{\Lambda}$. Put $\mathscr{X}$ = \{im$(f_{ji})$, im$(h_{i})$ \ $|$\  $i\geq 0, j\in J $\}}.\end{nota}

Now  we give a characterization of modules of virtually finite injective dimension   over right
Noetherian rings.



\begin{thm}\label{thm1}  Let  $\Lambda$ be a right  Noetherian ring and keep the notation as above. Then
$(^\perp(\mathscr{X}^\perp)$, $\mathscr{X}^\perp)$ is a hereditary complete
  cotorsion theory. Moreover,  $^\perp\mathscr{GI}$ = $^\perp(\mathscr{X}^\perp)$, i.e.,
 $\mathscr{X}^\perp$ = $\mathscr{GI}$. \end{thm}

\begin{proof}
Since $\mathscr{X}$ is a set, $(^\perp(\mathscr{X}^\perp)$, $\mathscr{X}^\perp)$ is a   complete
  cotorsion theory by  \cite[Theorem 7.4.1]{eno5}. Next we prove that it is hereditary.

  Since  each im$(f_{ji})$ belongs to $^\perp(\mathscr{X}^\perp)$, for any $n \geq 1$, we deduce from the exact sequence $(\spadesuit_{j})$
that ${\rm Ext}_{\Lambda}^{n}($im$(f_{ji})$, $Y)=0$ for all $Y \in \mathscr{X}^\perp$. Therefore, im$(f_{ji}) \in$  $^{\perp_{\infty}}(\mathscr{X}^\perp)$.
On the other hand, every injective right  $\Lambda$-module is the direct sum
of copies of the various $E_{j}$. Hence   ${\bf Inj}(\Lambda) \subseteq $  $^{\perp_{\infty}}(\mathscr{X}^\perp)$.
But then we see from from the exact sequence $(\clubsuit)$
that  im$(h_{i}) \in$  $^{\perp_{\infty}}(\mathscr{X}^\perp)$.
   In addition, the class  $^\perp(\mathscr{X}^\perp)$ consists of all direct summands of $\mathscr{X}$-filtered modules by Lemma \ref{lem2}.
These observations  together with Lemma \ref{lem1} tell us that, ${\rm Ext}_{\Lambda}^{n}($$B$, $Y)=0$ for all $B \in {^\perp(\mathscr{X}^\perp)}$, $Y \in \mathscr{X}^\perp$,  and $n \geq 1$.
So the pair $(^\perp(\mathscr{X}^\perp)$, $\mathscr{X}^\perp)$   is hereditary by \cite[Lemma 5.24]{gbe}.

Now we prove that  $\mathscr{X}^\perp$ = $\mathscr{GI}$.
Since injective modules are contained in the thick class $^\perp\mathscr{GI}$ (cf. \cite[Lemma 5.24]{ss}), we see from the exact sequences $(\spadesuit_{j})$  and $(\clubsuit)$
that $\mathscr{X} \subseteq {^\perp\mathscr{GI}}$. Hence $\mathscr{X}^\perp \supseteq {(^\perp\mathscr{GI}})^\perp$ = $\mathscr{GI}$ (cf. Proposition \ref{prop1}).

It remains to show that $\mathscr{X}^\perp $   $\subseteq \mathscr{GI}$. For any $M \in \mathscr{X}^\perp $, there is an exact sequence
$0\rightarrow  K    \rightarrow \Lambda^{(\lambda)} \rightarrow  M\rightarrow 0   $  for some set $\lambda$. Then the sequence
$0\rightarrow  \Lambda^{(\lambda)}  \rightarrow E_{0}^{(\lambda)} \rightarrow {\rm(im(}h_{1}))^{(\lambda)}\rightarrow 0   $  is also exact.
 Construct a pushout of  $\Lambda^{(\lambda)}  \rightarrow E_{0}^{(\lambda)}$ along $\Lambda^{(\lambda)} \rightarrow  M$, we obtain a commutative diagram with exact
rows  and columns:


 $$\xymatrix{&&0\ar[d]&0\ar[d]&\\0\ar[r]& K\ar[r]
\ar@{=}[d]&\Lambda^{(\lambda)}\ar[d]\ar[r]
&M\ar[d]\ar[r]&0\\
0\ar[r]&  K\ar[r]
& E_{0}^{(\lambda)}\ar[d]\ar[r]
& Q\ar[d]\ar[r]&0\\
 & &{\rm im(}h_{1})^{(\lambda)}\ar[d]\ar@{=}[r]&{\rm im(}h_{1})^{(\lambda)}\ar[d]&&&\\
&&0&0}$$
Note that ${\rm Ext}_{\Lambda}^{1}($${\rm im(}h_{1})^{(\lambda)}$, $M)   \cong  $ $({\rm Ext}_{\Lambda}^{1}$${\rm(im(}h_{1})$, $M))^{\lambda} =  0$.
Thus  $Q \cong M \oplus {\rm im(}h_{1})^{(\lambda)}$, and we get   an epimorphism $E_{0}^{(\lambda)} \rightarrow  M$ with $E_{0}^{(\lambda)}  $ injective ($\Lambda$ is right  Noetherian).
Hence there exists a short exact sequence   $$0 \longrightarrow K \longrightarrow I_{0}  \stackrel {g_{0}}\longrightarrow M \longrightarrow 0$$
where  $g_{0}$ is an  injective cover.  Then $K \in {\bf Inj}(\Lambda)^\perp$ by \cite[Corollary 7.2.3]{eno5}. But   $M \in {\bf Inj}(\Lambda)^{\perp_{\infty}}$.
 Thus, one can see from the above short exact sequence that $K \in  {\bf Inj}(\Lambda)^{\perp_{\infty}}$. For any $i \geq 1$, we have an exact sequence
  $$0 = {\rm Ext}_{\Lambda}^{1}{\rm(im}(h_{i}), M)  \rightarrow  {\rm Ext}_{\Lambda}^{2}{\rm(im}(h_{i}), K)  \rightarrow  {\rm Ext}_{\Lambda}^{2}{\rm(im}(h_{i}), I_{0})  = 0.$$
So  $  {\rm Ext}_{\Lambda}^{2}{\rm(im}(h_{i}), K)   = 0$.  Now consider the short exact sequence
$0\rightarrow  {\rm im(}h_{i - 1})  \rightarrow E_{i - 1}\rightarrow {\rm im(}h_{i})\rightarrow 0   $
which  induces an exact sequence  $$0 = {\rm Ext}_{\Lambda}^{1}(E_{i - 1}, K)  \rightarrow  {\rm Ext}_{\Lambda}^{1}{\rm(im}(h_{i-1}), K)  \rightarrow  {\rm Ext}_{\Lambda}^{2}{\rm(im}(h_{i}), K)  = 0.$$
Thus $  {\rm Ext}_{\Lambda}^{1}{\rm(im}(h_{i-1}), K)   = 0$ for any $i \geq 1$. This implies  $K \in \mathscr{X}^{\perp}$.
Recursively, we construct an exact sequence of right $\Lambda$-modules
$$ \cdots \stackrel {g_{2}} \longrightarrow
I_1\stackrel {g_{1}} \longrightarrow I_0\stackrel {g_{0}}
\longrightarrow M  \longrightarrow  0  $$
with $I_i \in {\bf Inj}(\Lambda)$ and ker$(g_i) \in {\bf Inj}(\Lambda)^{\perp_{\infty}}$ ($i \geq 0$).
Therefore, $M $   $\in \mathscr{GI}$ by \cite[Lemma 17(2)]{gii}, and   $\mathscr{X}^\perp $   $\subseteq \mathscr{GI}$, as required. \end{proof}



\medskip


\section{Proof of the Gorenstein Symmetry Conjecture for artin algebras }


We begin with the following definition.

\begin{defn} \label{defn3} \rm{ Let $M$ be a    $\Lambda$-module.  We say that  $M$ has {\it Gorenstein projective
dimension} at most  $n$, and we write  Gpd$_{\Lambda}$($M$) $\leq n$, if there exists an exact sequence
of  $\Lambda$-modules $ 0 \rightarrow G_n\rightarrow
\cdots \rightarrow G_0\rightarrow M\rightarrow 0$ where each $G_{i}$ is Gorenstein projective. If there is no such $n$, set   Gpd$_{\Lambda}$($M$)
$=\infty$.

The {\it Gorenstein injective
dimension},  Gid$_{\Lambda}$($M$), is     defined similarly.}\end{defn}

Bennis   and Mahdou    \cite[Theorem 1.1]{bem} proved the     following equality
\\  \\
\centerline{   sup\{Gpd$_{\Lambda}(M)$\ $|$\ $M$     is  a  right  $\Lambda$-module\} =  sup\{Gid$_{\Lambda}(M)$\ $|$\ $M$     is  a  right  $\Lambda$-module\}.}\\ \\
The  common value of the above quantities is called the {\it right Gorenstein global
dimension} of $\Lambda$ and we denote it by     $r$.Ggldim($\Lambda$).

It is  known that $   r.{\rm Ggldim(}\Lambda)$  is related to the cohomological invariants $r$.silp($\Lambda$) := sup\{id$_{\Lambda}(P)$\ $|$\ $P\in$     ${\bf Proj}(\Lambda)$\} and $r$.spli($\Lambda$)  := sup\{pd$_{\Lambda}(E)$\ $|$\ $E \in$    ${\bf Inj}(\Lambda)$\}
introduced by Gedrich and Gruenberg in \cite{ged}  via the formula  \\
 \\
\centerline{ {\rm  $   r.{\rm Ggldim(}\Lambda)$ }{\rm = max\{$r$.silp($\Lambda$),  $r$.spli($\Lambda$)\}} \ \ \ (\cite[Theorem 3.3]{bou1} and  \cite[Theorem 2.1]{mad1}).}
 \\
 \\
In addition,  if  $   r.{\rm Ggldim(}\Lambda) < \infty$, then {\rm  $   r.{\rm Ggldim(}\Lambda)$ } = $r$.silp($\Lambda$) =  $r$.spli($\Lambda$) by   \cite[Proposition 2.3]{mad1}.

As usual, we call a ring $\Lambda$  Noetherian in case it is both left Noetherian and right Noetherian.

\begin{prop}\label{prop2}  Let $\Lambda$  be    a  Noetherian ring. Then \\
 \\
\centerline{   $   r.{\rm Ggldim(}\Lambda)$  =  $   r.{\rm Ggldim(}\Lambda^{{\rm op}})$.  } \end{prop}

\begin{proof} This follows from  \cite[Corollary 2.7]{lyz}; see also  \cite[Theorem 3.10]{bou1}. \end{proof}

 Recall that an epimorphism   $f:$ $ A\twoheadrightarrow B$ of $\Lambda$-modules is said to be   {\it pure}  if the map   $f^{\ast}:$ Hom$_{\Lambda}$$(F, A)\rightarrow$ Hom$_{\Lambda}$$(F, B)$
is epic for all finitely presented   $\Lambda$-modules $F$; then ker($f$) is said to be  {\it pure} in $A$.

\begin{prop}\label{prop3}  Let $\Lambda$  be    a right Noetherian ring, $\kappa$ an infinite cardinal with {\rm Card($\Lambda$)}  $\leq \kappa$.
Let $\mathscr{GI}_{\kappa}$  denote  an irreduntant set of representatives of modules $G$  such that $G \in$ $\mathscr{GI}$ and {\rm Card($G$)}  $\leq \kappa$.
 Assume that the class   $\mathscr{GI}$   is closed under pure submodules.
Then  any  $M\in \mathscr{GI}$  is $\mathscr{GI}_{\kappa}$-filtered. \end{prop}

\begin{proof} First, we prove that $\mathscr{GI}$ is closed under  direct limits.

 Let \{$G_{i}\}_{i\in I}$ be any direct system of   modules from  $\mathscr{GI}$. Then we get a pure exact sequence
$0\rightarrow A\rightarrow \bigoplus G_{i} \rightarrow \lim\limits_{ \longrightarrow }G_{i} \rightarrow 0$   by  \cite[Corollary 2.9]{gbe}.  Note that $\bigoplus G_{i} $ is pure in $\prod G_{i} $ by \cite[Corollary 2.21(c)]{gbe}. So both $\bigoplus G_{i} $ and  $A$ belong to $ \mathscr{GI}$.   It follows that  $\lim\limits_{ \longrightarrow }M_{i}$ $\in \mathscr{GI}$ by Proposition   \ref{prop1}.

Let $M \in$ $\mathscr{GI}$. We mimic  the proof     of \cite[Corollary 3.3.4]{eno6} to show that $M$  is $\mathscr{GI}_{\kappa}$-filtered. For any $x\in$ $M$, by  \cite[Lemma
5.3.12]{eno5}, there is a pure submodule  $M_{0}$ of $M$ with $x\in$ $M_{0}$ such that
Card($M_{0}$) $\leq \kappa$	 (simply let  $N = x\Lambda$ and $f $ the   inclusion map from $N $ to $M $ in the lemma). We see that both  $M_{0}$ and  $M/M_{0}$
are in $\mathscr{GI}$ by Proposition \ref{prop1}.

For any  $x_{1} + M_{0} \in$ $M/M_{0}$, again by  \cite[Lemma
5.3.12]{eno5}, there is a pure submodule  $M_{1}/M_{0}$ of $M/M_{0}$ such that $x_{1} +  M_{0}  \in$ $M_{1}/M_{0}$ and
Card($M_{1}/M_{0}$) $\leq \kappa$.  But  $M_{1}/M_{0} \in$ $\mathscr{GI}$,  so it is isomorphic to an element of $\mathscr{GI}_{\kappa}$.
Note that  $\mathscr{GI}$   is closed under extensions. Hence $M_{1} \in$ $\mathscr{GI}$.
  Thus   $M/M_{1} \in$ $\mathscr{GI}$ again by Proposition \ref{prop1}.

  Proceeding by induction, we get a chain of  modules $M_{n} \in$ $\mathscr{GI}$,  $n \in$ $\mathbb{N}$, such that $M_{n+1}/M_{n}  $ is isomorphic to an element of $\mathscr{GI}_{\kappa}$.
Let $M_{\omega_{0}}$ = $\bigcup_{n\in \mathbb{N}}M_{n}$, then $M_{\omega_{0}}$ $\in \mathscr{GI}$ and  $M/M_{\omega_{0}}$ $\in \mathscr{GI}$.
Choosing $x_{\omega_{0}+1} + M_{\omega_{0}} \in M/M_{\omega_{0}}$, we find  a pure submodule  $M_{\omega_{0}+1}/M_{\omega_{0}}$ of $M/M_{\omega_{0}}$ such that
 $x_{\omega_{0}+1} + M_{\omega_{0}}  \in$ $M_{\omega_{0}+1}/M_{\omega_{0}}$,
Card($M_{\omega_{0}+1}/M_{\omega_{0}}$) $\leq \kappa$ and  $M_{\omega_{0}+1}/M_{\omega_{0}} \in$ $\mathscr{GI}$.  So $M_{\omega_{0}+1}/M_{\omega_{0}}$  is isomorphic to an element of $\mathscr{GI}_{\kappa}$.

We have shown that  $\mathscr{GI}$  is closed under direct limits. Hence, by transfinite induction  we can write
$M$ as a union of a continuous chain $(M_{\alpha})$$_{\alpha < \lambda}$ of  pure submodules of $M$, such that
 $M_{0}\in$  $\mathscr{GI}_{\kappa}$, $M_{\alpha + 1}/M_{\alpha}  $  is isomorphic to an element of $\mathscr{GI}_{\kappa}$  whenever $\alpha + 1 < \lambda$,
i.e., $M$  is $\mathscr{GI}_{\kappa}$-filtered. The
proof is finished.\end{proof}

To simplify the proof of the next theorem, we give the following notation.
\begin{nota}\label{nota2}  {\rm  \  Let  $\Lambda$ be a  right Noetherian ring and ($E_{j})_{j\in J}$ the family be as in Notation \ref{nota1}.
Let  $\kappa$ be an infinite cardinal with {\rm Card($\kappa$)}   $\geq$ max\{Card($\Lambda$),  Card($\bigoplus_{j\in J}E_{j})$\}.
Let $\mathscr{U}$  denote  an irreduntant set of representatives of  right  $\Lambda$-modules $A$  such that Gid$_{\Lambda}$($A$)  $< \infty$ and {\rm Card($A$)}  $\leq \kappa$.
 For each $A_{u} \in \mathscr{U}$, choose  a  projective resolution $$ \cdots \stackrel {\pi_{u2}} \longrightarrow
{P_{u1}}\stackrel {\pi_{u1}} \longrightarrow {P_{u0}}\stackrel {\pi_{u0}}
\longrightarrow A_{u}  \longrightarrow  0  \eqno (\heartsuit_{u})$$
 of  $A_{u}$, and let
 $$ 0  \longrightarrow A_{u}\stackrel {\tau_{u0}}
\longrightarrow
I_{u0}\stackrel  {\tau_{u1}} \longrightarrow I_{u1}\stackrel  {\tau_{u2}}
\longrightarrow \cdots \eqno (\sharp_{u})$$
  be a minimal  injective resolution of   $A_{u}$. Put $\overline{\mathscr{U}}$ = \{im$(\pi_{ui})$, im$(\tau_{ui})$ \ $|$\  $i\geq 0, u\in \mathscr{U} $\}}.\end{nota}

 We believe that the following result is of   interest.
\begin{thm}\label{thm2}  Let  $\Lambda$ be a right  Noetherian ring and keep the notation as above. If {\rm Gid}$_{\Lambda}$($\Lambda$)  $< \infty$ and the class   $\mathscr{GI}$   is closed under pure submodules,
then:
\begin{itemize}
\item[$(1)$]  every right  $\Lambda$-module  $M$ is   isomorphic to a  direct summand a    $\overline{\mathscr{U}}$-filtered module;
 \item[$(2)$] {\rm Gid}$_{\Lambda}(S)  < \infty$ for any simple module  $S$. \end{itemize}\end{thm}

\begin{proof} (1) \ Since  {\rm Gid}$_{\Lambda}$($\Lambda$)  $< \infty$ and   {\rm Card($\kappa$)} $\geq$ max\{Card($\Lambda$),  Card($\bigoplus_{j\in J}E_{j})$\},
 we may assume that both $\Lambda$ and  $\bigoplus_{j\in J}E_{j}$ belong to the set $\mathscr{U}$.
Hence  we may assume that  the set $\mathscr{X}$ constructed in Notation  \ref{nota1}  is contained in $\overline{\mathscr{U}}$.  Therefore, $ \overline{\mathscr{U}}^\perp $  $\subseteq \mathscr{X}^\perp$ =
  $\mathscr{GI}$ (cf. Theorem \ref{thm1}),
 and $(^\perp(\overline{\mathscr{U}}^\perp)$, $\overline{\mathscr{U}}^\perp)$ is a   complete
  cotorsion theory by  \cite[Theorem 7.4.1]{eno5}. Next we prove that the pair $(^\perp(\overline{\mathscr{U}}^\perp)$, $\overline{\mathscr{U}}^\perp)$ is hereditary.

  Since  each im$(\pi_{ui})$ belongs to $^\perp(\overline{\mathscr{U}}^\perp)$, for any $n \geq 1$, we deduce from the exact sequence $(\heartsuit_{u})$
that ${\rm Ext}_{\Lambda}^{n}($im$(\pi_{ui})$, $Y)=0$ for all $Y \in \overline{\mathscr{U}}^\perp$. Therefore, im$(\pi_{ui}) \in$  $^{\perp_{\infty}}(\overline{\mathscr{U}}^\perp)$.
On the other hand, every injective right  $\Lambda$-module is the direct sum
of copies of the various $E_{j}$. Hence   ${\bf Inj}(\Lambda) \subseteq $  $^{\perp_{\infty}}(\overline{\mathscr{U}}^\perp)$.
But then we see from from the exact sequence $(\sharp_{u})$
that  im$(\tau_{ui}) \in$  $^{\perp_{\infty}}(\overline{\mathscr{U}}^\perp)$.
   In addition, the class  $^\perp(\overline{\mathscr{U}}^\perp)$ consists of all direct summands of $\overline{\mathscr{U}}$-filtered modules by Lemma \ref{lem2}.
These observations  together with Lemma \ref{lem1} tell us that, ${\rm Ext}_{\Lambda}^{n}($$W$, $Y)=0$ for all $W \in {^\perp(\overline{\mathscr{U}}^\perp)}$, $Y \in \overline{\mathscr{U}}^\perp$,  and $n \geq 1$.
So the pair $(^\perp(\overline{\mathscr{U}}^\perp)$, $\overline{\mathscr{U}}^\perp)$   is hereditary by \cite[Lemma 5.24]{gbe}.

  For any $M \in $ Mod-$\Lambda$, there is an exact sequence
$0\rightarrow  M    \rightarrow Q \rightarrow  L\rightarrow 0   $ with  $Q\in $ $\overline{\mathscr{U}}^\perp$ $\subseteq \mathscr{GI}$ and  $L \in $ $^\perp(\overline{\mathscr{U}}^\perp)$.
Note that the class   $\mathscr{GI}$   is closed under pure submodules. So $Q $ is $\mathscr{GI}_{\kappa}$-filtered by Proposition \ref{prop3}.
By  construction,  every module in $\mathscr{GI}_{\kappa}$  is isomorphic to an element of $\overline{\mathscr{U}}$.  Hence $Q $ is $\overline{\mathscr{U}}$-filtered.
This implies $Q \in $ $^\perp(\overline{\mathscr{U}}^\perp)$. It follows that $M\in $ $^\perp(\overline{\mathscr{U}}^\perp)$. This proves (1).


 Now let   $S\in $ Mod-$\Lambda$ be simple. Since $S\in $ $^\perp(\overline{\mathscr{U}}^\perp)$, $S$ is  isomorphic to a  direct summand of some module $Y$  in
 $\overline{\mathscr{U}}$ by Corollary  \ref{cor1}. According to the construction of $\overline{\mathscr{U}}$, {\rm Gid}$_{\Lambda}$($Y$)  $< \infty$.
Then the  the dual   of \cite[Proposition 2.19]{hol} forces that {\rm Gid}$_{\Lambda}$($S$)  $< \infty$.  This   finishes  the proof.
  \end{proof}

   The next lemma is well-known, but we could not find a concrete reference to it.
  \begin{lem}\label{lem3}  Let
$0\rightarrow  B    \rightarrow C\rightarrow  D\rightarrow 0   $ be a  short exact sequence in {\rm Mod-}$\Lambda$.
If  {\rm Gid}$_{\Lambda}$($B$)  $\leq n < \infty$ and   {\rm Gid}$_{\Lambda}$($D$)  $\leq n$, then  {\rm Gid}$_{\Lambda}$($C$)  $\leq n $.\end{lem}
\begin{proof} This is a direct consequence of  {\rm \cite[Theorems  2.22 and 2.24]{hol}}. \end{proof}

 We also need the following   result.
\begin{lem} \label{lem4} Let  $n\geq 0$  be given and   let $\mathscr{GI}_{n}$ denote  the class of right $\Lambda$-modules $M$
 such that   {\rm Gid}$_{\Lambda}${\rm (}$M${\rm )}  $\leq n.$  If $\mathscr{GI}$ is closed under direct limits, then $\mathscr{GI}_{n}$ is also closed under direct limits.\end{lem}

\begin{proof} Let \{$M_{i}\}_{i\in I}$ be any direct system of
right $\Lambda$-modules   with each  $M_{i} \in$  $\mathscr{GI}_{n}$.
 Take an  injective resolution
 $$ 0  \longrightarrow M_{i}\stackrel {\theta_{i0}}
\longrightarrow
G_{i0}\stackrel {\theta_{i1}} \longrightarrow G_{i1}\stackrel {\theta_{i2}}
\longrightarrow \cdots   $$
 for each  $M_{i}$
and let $L_{in} =$ im($\theta_{in}$). Then   $L_{in} $  $\in \mathscr{GI}$ by \cite[Theorem  2.22 (iv)]{hol}. Since       $ \mathscr{GI}$
 is closed under direct limits, we see from the exactness  of the following complex
  $$0 \rightarrow \lim\limits_{ \longrightarrow }M_{i}\rightarrow \lim\limits_{ \longrightarrow }G_{i0}\rightarrow \lim\limits_{ \longrightarrow }G_{i1} \rightarrow
\cdots    \rightarrow \lim\limits_{ \longrightarrow }L_{in}\rightarrow 0$$
that  $\lim\limits_{ \longrightarrow }M_{i} $  $ \in$  $\mathscr{GI}_{n}$, as required.
\end{proof}



Now we turn our attention to artinian rings.

\begin{thm}\label{thm3}  Let  $\Lambda$ be an Artinian ring. If   the   classes of all Gorenstein     injective
   right and left  $\Lambda$-modules   are closed under pure submodules, then the following quantities are identical:
\begin{itemize}

\item[$(1)$] $r.{\rm Ggldim(}\Lambda)$;
\item[$(2)$] $r.{\rm Ggldim(}\Lambda^{{\rm op}})$;
\item[$(3)$]  {\rm id}$_{\Lambda}$($\Lambda$);
 \item[$(4)$] {\rm id}$_{\Lambda^{{\rm op}}}$($\Lambda^{{\rm op}}$);
 \item[$(5)$]  {\rm Gid}$_{\Lambda}$($\Lambda$);
 \item[$(6)$] {\rm Gid}$_{\Lambda^{{\rm op}}}$($\Lambda^{{\rm op}}$).\end{itemize}\end{thm}

 \begin{proof} It is clear that (5) $\leq (3)$, and \cite[Theorem 3.3]{bou1} tells us that  (3) $\leq (1)$.
 Next we prove that (1) $\leq (5)$. We may assume that   {\rm Gid}$_{\Lambda}$($\Lambda$)  $< \infty$.
Since $\Lambda$ is  Artinian,   there is an irredundant finite set   $\mathscr{S}$ of representatives of the simple modules in Mod-$\Lambda$.
Then  {\rm Gid}$_{\Lambda}(S)  < \infty$ for any $S\in $  $\mathscr{S}$ by Theorem \ref{thm2}(2).
 So $n$ = sup\{Gid$_{\Lambda}(S)$ \ $|$\   $S\in $  $\mathscr{S}$ \} $  < \infty$. In addition, it is well-known that any $M \in $ mod-$\Lambda$
 is  finitely  $\mathscr{S}$-filtered (cf. \cite[Proposition 11.1]{and}). By induction, we get from Lemma \ref{lem3} that
 Gid$_{\Lambda}(M) \leq n.$ But every   module   is the  direct limit of its  finitely generated submodules, and $\mathscr{GI}_{n}$ is closed under direct limits  by Lemma \ref{lem4}.
  Thus  Gid$_{\Lambda}(Y) \leq n$  for any    $Y\in $ Mod-$\Lambda$.  It follows from \cite[Theorem 1.1]{bem} that $r.{\rm Ggldim(}\Lambda) \leq n$. We then deduce from
 \cite[Lemma 2.1]{bem} that  $r.{\rm Ggldim(}\Lambda) =$  {\rm id}$_{\Lambda}$($\Lambda$) $< \infty$. But then we get  from   the dual of  \cite[Proposition 2.27]{hol} that {\rm id}$_{\Lambda}$($\Lambda$)
 = {\rm Gid}$_{\Lambda}$($\Lambda$).
Hence (1) $\leq (5)$. Therefor, (1) $= (3) = (5)$.

 The proof of  (2) $= (4) = (4)$  is symmetric.

Finally, apply Proposition \ref{prop2} to get that (1) $= (2) = (3)$   $= (4) = (5) = (6)$.
 \end{proof}

   We are now in a position to show the validity of  the Gorenstein Symmetry Conjecture
for all  Artin algebras.

\begin{thm}\label{thm4}  Let  $\Lambda$ be an Artin algebra. Then the following quantities are identical:
\begin{itemize}

\item[$(1)$] $r.{\rm Ggldim(}\Lambda)$;
\item[$(2)$] $r.{\rm Ggldim(}\Lambda^{{\rm op}})$;
\item[$(3)$]  {\rm id}$_{\Lambda}$($\Lambda$);
 \item[$(4)$] {\rm id}$_{\Lambda^{{\rm op}}}$($\Lambda^{{\rm op}}$);
 \item[$(5)$]  {\rm Gid}$_{\Lambda}$($\Lambda$);
 \item[$(6)$] {\rm Gid}$_{\Lambda^{{\rm op}}}$($\Lambda^{{\rm op}}$).\end{itemize}\end{thm}

 \begin{proof} Since $\Lambda$ is an Artin algebra, the   classes of all Gorenstein     injective
   right and left  $\Lambda$-modules   are closed under pure submodules by \cite[Proposition 3.8(iii)]{bel}. Now apply Theorem \ref{thm3}.
 \end{proof}


\medskip




   \vskip3mm

\begin{thebibliography}{99}

\bibitem{and} F. W. Anderson,  K. R. Fuller,  {\it Rings and
Categories  of  Modules}, 2nd. Springer-Verlag, New York, 1992.

\bibitem{ar1}  M. Auslander, I. Reiten, Applications of contravariantly finite subcategories. {\it Adv. Math.} 86 (1991), 111-152.

\bibitem{ars} M. Auslander, I. Reiten, S. O. Smal${\o}$, {\it Representation Theory of Artin Algebras.} Cambridge Studies in Advanced Mathematics, vol. 36. Cambridge Univ.  Press, Cambridge, 1995.

\bibitem{bel} A. Beligiannis, Cohen-Macaulay modules, (co)torsion pairs and virtually Gorenstein algebras. {\it J. Algebra }  288 (2005), 137-211.

\bibitem{bou1} S. Bouchiba,   Finiteness aspects of Gorenstein homological dimensions. {\it Math. Colloq.} 131(2),  (2013), 171-194.  

\bibitem{bem} D. Bennis, N. Mahdou, Global Gorenstein dimensions. {\it Proc.
Amer.  Math.  Soc}.  138(2) (2010), 461-465.

\bibitem{eno5} E. E. Enochs, O. M. G. Jend, {\it Relative  Homological  Algebra},  Berlin,  Walter de Gruyter,
2000.

\bibitem{eno9}  E. E. Enochs, O. M. G. Jend, Gorenstein injective and projective modules,  {\it Math. Z}. 220(4)  (1995), 611-633. 

\bibitem{eno6} E. E. Enochs, L. Oyonarte, Covers, Envelopes and Cotorsion Theories, New York,
Nova Science Publishers, Inc., 2002.

 \bibitem{gii} J. Gillespie, A. Iacob, Duality pairs, generalized Gorenstein modules, and Ding injective
envelopes, {\it Comptes Rendus Math$\acute{{\rm e}}$matique},  360 (2022), 381-398. 

\bibitem{ged} T. V. Gedrich,  K. W. Gruenberg,  Complete cohomological functor on groups.
{\it Topology Appl.} 25 (1987), 203-223.  

\bibitem{hua} Z. Y. Huang, Selforthogonal modules with finite injective dimension II, {\it J. Algebra} 264 (2003),   262-268.

\bibitem{hol} H. Holm,  Gorenstein homological dimensions. {\it J. Pure Appl. Algebra} 189 (2004), 167-193.  

\bibitem{lyz} W. Q. Li, L. Yan,  D. D. Zhang,  Some results on Gorenstein (weak) global dimension of rings, {\it Comm. Algebra}, 51(1) (2023),  264-275.

\bibitem{mao}  L. X. Mao,  N. Q. Ding,  Relative cotorsion modules and relative flat modules, {\it Comm. Algebra}, 34 (2006), 2303-2317.

\bibitem{mad1} N.  Mahdou, M. Tamekkante, On (weak) Gorenstein global dimensions. {\it Acta Math. Univ. Comenianae} 2  (2013), 285-296.

\bibitem{gbe} R. G$\ddot{\rm o}$bel, J. Trlifaj, {\it Approximations and Endomorphism Algebras of Modules}, de Gruyter Exp. Math. vol.
41,  2 revised and extended edn. de Gruyter, Berlin, (2012).


 \bibitem{ss} J. $\check{\rm S}$aroch, J. $\check{\rm S}$$\check{\rm t}$ov$\acute{\rm {\i}}$$\check{\rm c}$ek,   Singular compactness and definability for $\sum$-cotorsion and Gorenstein modules,
{\it Sel. Math. New Ser.} 26:23 (2020). 
 \end{thebibliography}






\end{document}
